\documentclass[10pt]{article}
\usepackage[T1]{fontenc}
\usepackage[utf8]{inputenc}
\usepackage[french]{babel}
\usepackage{amsmath,amssymb,mathrsfs,mathpazo}
\usepackage{geometry}
\geometry{a4paper,margin=2cm}
\usepackage{graphicx,array,tabularx,tabularray,colortbl,multirow,multicol,booktabs,enumitem,pifont,ulem,fancybox,diagbox,cancel}
\usepackage{pstricks,pstricks-add,pst-plot,pst-tree,pst-node,pst-text,pst-func,pst-math,pst-bspline,pst-3dplot,pst-eucl,pst-solides3d}
\let\pscalculate\undefined
\usepackage{pst-bezier}
\usepackage{tikz,pgfplots,tkz-tab}
\usetikzlibrary{arrows,arrows.meta,patterns,calc,positioning,trees,decorations.markings,shapes.geometric,shapes.misc,shapes.arrows,intersections,decorations.pathmorphing,decorations.pathreplacing,backgrounds,shadows,babel}
\pgfplotsset{compat=1.18}
\usepgfplotslibrary{fillbetween}
\usepackage[np]{numprint}
\usepackage{listings,makecell,pas-tableur,twoopt}
\let\barre\undefined
\usepackage[pstricks]{tablvar}
\definecolor{darkgreen}{RGB}{0,100,0}
\providecommand{\red}{\color{red}}
\providecommand{\blue}{\color{blue}}
\providecommand{\green}{\color{green}}
\providecommand{\black}{\color{black}}
\providecommand{\white}{\color{white}}
\providecommand{\gray}{\color{gray}}
\providecommand{\cyan}{\color{cyan}}
\providecommand{\magenta}{\color{magenta}}
\providecommand{\yellow}{\color{yellow}}
\providecommand{\violet}{\color{violet}}
\providecommand{\orange}{\color{orange}}
\providecommand{\euro}{\texteuro}
\providecommand{\EUR}{\texteuro}
\providecommand{\ieme}{\textsuperscript{e}}
\providecommand{\degres}{\ensuremath{^\circ}}
\providecommand{\degree}{\ensuremath{^\circ}}
\providecommand{\vv}[1]{\overrightarrow{#1}}
\providecommand{\Coord}[1]{#1}
\providecommand{\decosix}{}
\providecommand{\starredbullet}{\ensuremath{\bullet}}
\pagestyle{empty}
\setlength{\parindent}{0pt}
\usepackage{adjustbox}\DeclareUnicodeCharacter{2236}{:}\DeclareUnicodeCharacter{2190}{\ensuremath{\leftarrow}}
\providecommand{\R}{\mathbb{R}}
\providecommand{\N}{\mathbb{N}}
\providecommand{\Z}{\mathbb{Z}}
\providecommand{\Q}{\mathbb{Q}}
\providecommand{\D}{\mathbb{D}}
\providecommand{\C}{\mathbb{C}}
\providecommand{\e}{{\mathrm{e}}}
\providecommand{\vect}[1]{\overrightarrow{#1}}
\providecommand{\vectt}[1]{\overrightarrow{\mathrm{#1}}}
\providecommand{\vvt}[1]{\overrightarrow{\mathrm{#1}}}
\providecommand{\np}[2][]{{#2~#1}}
\providecommand{\vv}[1]{{\overrightarrow{#1}}}
\renewcommand{\d}{{\mathrm{d}}}
\providecommand{\degres}{{^\circ}}
\providecommand{\degre}{{^\circ}}
\providecommand{\blue}{\color{blue}}
\providecommand{\red}{\color{red}}
\providecommand{\enskip}{\quad}
\providecommand{\llbracket}{\lbrack\!\lbrack}
\providecommand{\rrbracket}{\rbrack\!\rbrack}
\providecommand{\euro}{{\text{€}}}
\providecommand{\no}{{\text{n°}}}
\providecommand{\interval}[2]{{[#1;#2]}}
\providecommand{\intervFO}[2]{{[#1;#2[}}
\providecommand{\dotfill}{\ldots}

\providecommand{\ds}{\displaystyle}
\providecommand{\textding}[1]{#1}
\providecommand{\Oij}{\(\mathrm{O};\vec{\imath},\vec{\jmath}\)}
\providecommand{\Oijk}{\(\mathrm{O};\vec{\imath},\vec{\jmath},\vec{k}\)}
\providecommand{\Ouv}{\(\mathrm{O};\vec{u},\vec{v}\)}
\providecommand{\intervFF}[2]{\ensuremath{[#1;#2]}}
\renewcommand{\euro}{\ifmmode\text{\texteuro}\else\texteuro\fi}
\DeclareUnicodeCharacter{00B0}{\ifmmode^{\circ}\else\textdegree\fi}
\DeclareUnicodeCharacter{20AC}{\ifmmode\text{\texteuro}\else\texteuro\fi}
% Énoncés : Éducation nationale. Transcription et corrigés : APMEP, Denis Vergès et contributeurs. Découpage et classement : CoopMaths, Sébastien Lozano, Éric Elter et contributeurs. Adaptation de la présentation : Maths Explorer.
% CC BY-SA : https://coopmaths.fr/www/about/
% https://coopmaths.fr/alea/static/eam/2026/tex/eam_specifique_2026_06_asie_1_cor.tex
\begin{document}
\noindent\textit{APMEP / CoopMaths. CC BY-SA. Présentation adaptée par Maths Explorer.}\par\medskip
\textbf{Affirmation 1 : Fausse.}

	Le budget total pour un vélo est égal à $300+45+60 = 405$~\euro (document n°1), soit plus de 400~\euro.

	Donc, pour cinq vélos, le montant serait \np[\euro]{2025}, supérieur à \np{2000}~\euro, qui est le montant maximum prévu par la proviseure (document n°4).

	On ne pourra pas installer cinq vélos avec ce budget, l'affirmation est fausse.

\vspace{5mm}
\textbf{Affirmation 2 : Vraie.}

	En pédalant pendant une heure (60 minutes), on récupère \np[~kJ]{100} (document n°1).

	Il y a une relation de proportionnalité entre la durée du pédalage et l'énergie récupérée.

	Pour charger un téléphone portable, il faut \np[~kJ]{40} (document n°2):

	\begin{tabular}{c|c|c}
		durée&60  &$x$ \\\hline
		énergie&100  & 40 \\
	\end{tabular}\quad
	Il faut $\dfrac{40\times60}{100} = \dfrac{\np{2400}}{100} = 24$ minutes de pédalage pour récupérer \np[~kJ]{40}, et ainsi pouvoir charger un téléphone portable, donc l'affirmation est vraie.

\vspace{5mm}
\textbf{Affirmation 3 : Vraie.}

	Comme on choisit un élève au hasard, on est en situation d'équiprobabilité, donc les proportions sont assimilables à des probabilités.

	Parmi les 560 élèves favorables au projet, 180 sont des élèves de première (document n°3), donc la probabilité recherchée est $\dfrac{180}{560} = \dfrac{18}{56} = \dfrac{9}{28}$.

	On se souvient que $7\times4 = 28$ : le quart de 28 est donc 7, et donc $\dfrac{9}{28}$ est plus grand qu'un quart, l'affirmation est vraie.

\vspace{5mm}
\textbf{Affirmation 4 :}

	Deux conditions doivent être remplies pour que le projet soit lancé (document n°4):
	\begin{itemize}
		\item plus de la moitié des élèves du lycée doivent se déclarer favorables à cette initiative : c'est le cas, car le nombre d'élèves favorables est 560, ce qui est plus de la moitié des \np{1000} élèves du lycée (document n°3);

		\item au moins 60~\% des élèves de première doivent également y être favorables : c'est aussi le cas puisqu'il y a 300 élèves de première, dont 180 soit exactement 60~\% sont favorables au projet (document n°3), en effet, 60\,\% de 100 élèves, c'est 60 élèves, donc 60\,\% de 300 élèves, c'est le triple, soit $60 \times 3 = 180$.
	\end{itemize}
	Les deux conditions sont remplies, le projet pourra être lancé et l'affirmation est vraie.

\needspace{10\baselineskip}

\end{document}
